% Architecture
% This section documents the system architecture, module design,
% and the common interface shared by all storage engines.

\subsection{Overview}

The SFDB system is organised into three layers:

\begin{enumerate}
    \item \textbf{Canonical Layer}: The common semantic representation
        (Section~\ref{sec:semantic-model}) that all engines map to and
        from.
    \item \textbf{Engine Layer}: Three independent storage engines that
        share the same abstract interface.
    \item \textbf{Benchmark Layer}: A common evaluation framework that
        communicates only through the abstract interface, ensuring fair
        comparisons.
\end{enumerate}

Figure~\ref{fig:architecture} shows the component diagram.

\begin{figure}[ht]
\centering
\begin{tikzpicture}[node distance=2cm, every node/.style={draw, rounded corners, minimum width=4cm, align=center}]
    \node (query) {Query Language};
    \node (canonical) [below of=query] {Canonical Semantic Model};
    \node (kg) [below left=1.5cm and 0.5cm of canonical] {Knowledge Graph};
    \node (sheaf) [below of=canonical] {Sheaf Database};
    \node (simp) [below right=1.5cm and 0.5cm of canonical] {Simplicial Database};
    \node (bench) [below=2.5cm of canonical] {Benchmark Framework};

    \draw[->] (query) -- (canonical);
    \draw[->] (canonical) -- (kg);
    \draw[->] (canonical) -- (sheaf);
    \draw[->] (canonical) -- (simp);
    \draw[->] (kg) -- (bench);
    \draw[->] (sheaf) -- (bench);
    \draw[->] (simp) -- (bench);
\end{tikzpicture}
\caption{System architecture showing the three storage engines and
         their relationship to the canonical model and benchmark
         framework.}
\label{fig:architecture}
\end{figure}

\subsection{Module Design}

The system is decomposed into the following modules, each with a single
well-defined responsibility:

\begin{description}
    \item[Parser] Transforms a query string into a logical plan that is
        independent of any storage engine.
    \item[Validator] Checks that every \texttt{SemanticFact} satisfies
        its invariants---immutability, identifier uniqueness, confidence
        bounds, temporal consistency, and referential integrity.
    \item[Serializer] Converts \texttt{SemanticFact} instances to and
        from bytes in three formats: JSON, Parquet, and MessagePack.
        All serialisers are deterministic and lossless.
    \item[Importer / Exporter] Bulk data transfer between external
        formats (RDF/XML, Turtle, CSV, JSON-LD) and the canonical
        model.
    \item[Storage] The logical layer that maps canonical fields to
        engine-specific data structures and provides the CRUD interface.
    \item[Indexes] Engine-private auxiliary structures for fast lookup.
    \item[Query Planner] Transforms a logical query into a physical
        execution plan.  The same logical query produces different
        physical plans for each engine.
    \item[Query Optimizer] Applies cost-based optimisation rules to the
        physical plan using engine-specific cost models.
    \item[Execution Engine] Executes a physical plan against the storage
        engine and returns canonical results.
    \item[Statistics] Tracks cardinalities, index sizes, and operation
        counts for use by the optimizer and benchmark.
    \item[Visualization] Generates plots from benchmark results using
        \texttt{matplotlib}.
    \item[Benchmark Runner] Orchestrates comparative evaluation across
        all engines with metrics collection and semantic equivalence
        verification.
    \item[Paper Generator] Produces \LaTeX{} output from benchmark
        results for inclusion in this document.
\end{description}

\subsection{Common Interface}

Every storage engine implements the abstract \texttt{DatabaseEngine}
interface (Table~\ref{tab:common-interface}), guaranteeing that engines
are drop-in replaceable and that comparisons are fair.

\begin{table}[ht]
\centering
\caption{Common \texttt{DatabaseEngine} interface methods.}
\label{tab:common-interface}
\begin{tabular}{@{}ll@{}}
\toprule
Method & Description \\
\midrule
\texttt{create(config)} & Initialise storage and indexes \\
\texttt{drop()} & Destroy all data \\
\texttt{insert(fact)} & Store a single \texttt{SemanticFact} \\
\texttt{update(id, fact)} & Replace an existing fact \\
\texttt{delete(id)} & Remove a fact by identifier \\
\texttt{query(query)} & Execute a logical query \\
\texttt{explain(query)} & Return execution plan without running \\
\texttt{statistics()} & Return engine statistics \\
\texttt{verify()} & Run internal consistency checks \\
\texttt{export(fmt)} & Export all facts as bytes \\
\texttt{import(data, fmt)} & Import facts from bytes \\
\bottomrule
\end{tabular}
\end{table}

\subsection{Canonical Pipeline}

All data flows through the same canonical pipeline:

\begin{enumerate}
    \item A \texttt{SemanticFact} enters the system.
    \item The \textbf{Validator} checks invariants.
    \item The \textbf{Importer} or direct insertion sends the fact to
        each engine via \texttt{insert()}.
    \item Each engine's \textbf{Storage} layer maps the canonical fact
        to its internal representation.
    \item On query, the \textbf{Query Planner} and \textbf{Optimizer}
        produce an engine-specific physical plan.
    \item The \textbf{Execution Engine} runs the plan and returns
        \texttt{SemanticFact} instances.
    \item The \textbf{Benchmark Runner} collects metrics and verifies
        that results from all engines are semantically equivalent.
\end{enumerate}
